\section{Introduction}
\label{sec:introduction}
\subsection{Research background and related work}
Label distribution learning (LDL) represents each instance by a vector of non-negative description degrees summing to one~\cite{G1}. Unlike conventional multi-label annotations that merely indicate label presence or absence, LDL preserves rich quantitative information regarding the relative importance of different labels. LDL has demonstrated success across diverse domains, including crowd counting, facial expression recognition, and physiological prediction~\cite{ZW1,LS,QH}. However, converting native distributions into dominant labels or binary relevance sets inevitably loses fine-grained distributional information. Effective attribute reduction for LDL must therefore directly evaluate and respect the underlying distribution.

High-dimensional attribute spaces frequently contain redundant or uninformative features. While selecting an optimal feature subset enhances model interpretability and lowers measurement overhead, reducing attribute cardinality does not inherently guarantee improved predictive performance. Furthermore, simple cardinality metrics fail to account for deployment costs when individual attributes carry different acquisition burdens. Cost-sensitive learning addresses various operational expenses, such as test cost, misclassification cost, and computational or human resources~\cite{Turney2000}; for example, cost-sensitive embedded feature selection has been explored within support vector machines~\cite{AL}. In rough set theory, test-cost-sensitive reduction seeks optimal attribute subsets under explicit budget constraints: Min et al.~\cite{Min2014TestCost} formulated cost-constrained selection as a constraint satisfaction problem solved via backtracking heuristics, whereas decision-theoretic rough sets model attribute reliance using decision costs~\cite{Jia2013MinCost}. Although these studies share our concept of acquisition cost, they fundamentally rely on discrete decision classes. Existing LDL feature selectors—including granular and neighborhood models—typically score attributes based on information metrics and apply cardinality or threshold cuts. To the best of our knowledge, none formulate the task as finding a \emph{minimum-acquisition-cost subset subject to an explicit prediction-loss constraint}. This formulation shifts the conventional paradigm: rather than minimizing prediction loss for a fixed subset size, we minimize total acquisition cost while guaranteeing that a designated predictor's internal training loss remains within the full-attribute reference baseline. This directly addresses a critical deployment challenge: identifying the cheapest set of measurements that provably maintains predictive feasibility. Here, test cost specifically refers to the expense of acquiring attribute values for new instances during prediction, distinct from the offline execution time of the reduction algorithm itself.

Rough sets and neighborhood models offer a rigorous foundation for evaluating attribute-dependent information structures~\cite{CARS,NRS}. Extending these models to LDL requires adapting the supervisory signal from crisp classes to continuous distributions. Qian et al.~\cite{QX} incorporated feature similarity alongside label correlation, Deng et al.~\cite{DL} proposed a dual-similarity neighborhood fuzzy entropy spanning both attribute and label spaces, and Qian et al.~\cite{QD} handled incomplete distributions via neighborhood-tolerance discrimination. Recent extensions further enrich this lineage: Lin et al.~\cite{LW} integrated sparse regularization, label correlation, and mutual information into feature weighting; Shu et al.~\cite{Shu2024Specific} jointly modeled common and label-specific features; Deng et al.~\cite{Deng2025LDNF} formulated neighborhood fuzzy rough sets for LDL; Zhao et al.~\cite{Zhao2026IWNRS} developed an information-weighted neighborhood model with a conditional discrimination index; and Chen et al.~\cite{Chen2025OnlineLDL} explored online selection driven by label importance and correlation. Additionally, information-theoretic measures—such as mutual information, fuzzy entropy~\cite{QHM,WH2,XQ}, neighborhood self-information~\cite{WH1,XJ,JZ}, and hierarchical variants~\cite{LH}—evaluate feature utility based on distributional information content. While these methods explore diverse representations and search strategies, their existence underscores the need to define precisely what an acquisition-cost-aware granular framework adds to the field.

A closely related domain is feature reduction for multi-label learning, as surveyed by Kashef et al.~\cite{KP}. Representative approaches include ReliefF-based multi-label weighting~\cite{RM1}, multivariate mutual-information filters~\cite{LK}, and fuzzy relative discernibility ranking~\cite{QX2}. Although multi-label techniques handle label correlations, they ultimately target binary relevance structures and remain unequipped for full description-degree distributions. Similarly, metaheuristic attribute reduction algorithms—spanning evolutionary, swarm, and local search strategies~\cite{HX,NG,ZS,AE,OM}—optimize classical dependency or information criteria. Such search mechanisms are orthogonal to our primary contribution and could, in principle, complement or substitute the constructive search employed here.

Granular-ball computing has also emerged as a powerful paradigm in LDL feature selection. Qian et al.~\cite{Qian2023GBLDL} developed a granular-ball fuzzy rough set incorporating distributional similarity and local consistency. Granular generation itself remains a pivotal design choice: GBG++ provides fast and stable generation for classification~\cite{Xie2024GBG}, efficient adaptive generation methods continue to evolve~\cite{Xia2023GBadaptive}, and unified granular-ball models connect Pawlak and neighborhood rough sets within a common theoretical framework~\cite{Xia2023GBRS}. The present work builds upon this granular foundation by introducing a source-dependent information geometry, a deterministic distribution-driven ball construction, and an explicit acquisition-cost objective. Rather than merely applying granular balls to LDL, our motivation lies in formalizing the mathematical and empirical interplay between supervised geometry, regenerated granules, and prediction-feasible cost minimization.

\subsection{Motivation and contributions}
Three interconnected challenges motivate the proposed framework. First, the informative value of an attribute varies across individual objects. When the local mean label distribution within an attribute-induced neighborhood deviates from the global label distribution, that attribute exhibits local distributional discrimination. We quantify this association using a regularized divergence weight. Because this weight depends intrinsically on the source object, the resulting dissimilarity measure is directed. Additionally, a raw geometric term preserves subtle distinctions made by non-constant attributes that carry zero information weight. Critically, this distributional association serves as an exploratory selection signal rather than a strict guarantee of downstream predictive performance.

Second, altering the attribute subset dynamically reshapes the granular topology. Although the directed dissimilarity exhibits monotonicity under attribute inclusion when dataset statistics and normalization parameters remain fixed, regenerated granular balls are not necessarily nested. Consequently, granular dependency does not monotonically increase with attribute inclusion. Mistaking distance monotonicity for a dependency theorem leads to unsound pruning and stopping heuristics. To address this, our framework explicitly accommodates signed dependency gains and defines the exact structural properties that hold under the training representation.

Third, granular representation quality and predictor validation loss assess fundamentally different metrics. A subset yielding a superior training representation does not inherently generalize better to unseen label distributions. We therefore formulate attribute reduction as minimizing additive acquisition costs subject to an explicit, training-internal prediction-loss constraint. Granular dependency guides the constructive search trajectory, while the target predictor's internal loss enforces feasibility. Figure~\ref{fig:intro-roles} illustrates this functional separation, which retains the structural benefits of granular computing without relying on unjustified predictive guarantees derived purely from dependency metrics.

\begin{figure}[htbp]
	\centering
	\begin{tikzpicture}[>=Stealth, every node/.style={font=\small}, box/.style={draw,rounded corners,align=center,text width=3.05cm,minimum height=1.15cm}]
		\node[box] (w) at (0,0) {Label-distribution\\weights and geometry};
		\node[box] (g) at (4,0) {Regenerated balls\\and dependency};
		\node[box] (h) at (4,-2.2) {Cost-sensitive\\candidate ranking};
		\node[box] (l) at (0,-2.2) {Inner-fold\\prediction loss};
		\node[box] (f) at (0,-4.4) {Full-set reference\\feasibility constraint};
		\node[box] (o) at (4,-4.4) {Selected subset\\and acquisition cost};
		\draw[->] (w)--(g);
		\draw[->] (g)--node[right]{guides}(h);
		\draw[->] (l)--(h);
		\draw[->] (l)--(f);
		\draw[->] (h)--(o);
		\draw[->] (f.south)--++(0,-0.55)-| (o.south);
		\node[font=\scriptsize] at (2,-5.5) {feasibility check};
	\end{tikzpicture}
	\caption{Decoupled roles of granular information and predictive feasibility. The schematic depicts information flow during attribute selection; internal constraint satisfaction does not guarantee loss preservation on held-out test sets.}
	\label{fig:intro-roles}
\end{figure}

The principal contributions of this work are summarized as follows:
\begin{enumerate}
	\item \textbf{LDL-Driven Directed Dissimilarity:} We construct a local information weight tailored to label distributions and combine it with raw attribute geometry to formulate a source-dependent directed dissimilarity. We prove its boundedness and fixed-statistics attribute monotonicity without requiring metric axioms or asserting monotonicity for regenerated granular dependency.
	\item \textbf{Deterministic Granular Dependency Framework:} We introduce distribution consistency and eligible-ball dependency within a deterministic granular ball construction. Through a within--between error decomposition, we show that dependency measures normalized training representation error reduction (using global-distribution fallback for ineligible objects), where signed dependency gains naturally capture structural variations under feature addition.
	\item \textbf{Constrained Acquisition-Cost Optimization:} We formalize minimum-acquisition-cost reduction under an internal prediction-loss constraint, proposing both a dependency-guided heuristic search and an exact backtracking reference algorithm. We establish the finite feasible termination of the heuristic and the global cost optimality of the completed exact search. Furthermore, an optional local deletion/exchange post-processing phase is designed to strictly lower cost while maintaining feasibility.
	\item \textbf{Rigorous Empirical Benchmarking:} We evaluate our framework across twenty-five native LDL benchmarks using repeated training partitions, reserved audit labels, matched-cost label-free controls, and budget-matched ablations, against five published scoring cores spanning 2022--2025. Results demonstrate that our method reduces acquisition cost by over two orders of magnitude (e.g., from $1648$ to approximately $10$ on the Natural Scene dataset) while consistently outperforming label-free baselines of equivalent cost. A guard ablation shows that conventional guard-free selection spends acquisition budget for predictions worse than the free default on the degenerate domains, and a predictor-sensitivity study shows the guard certifying informative low-cost subsets (dominating both zero-cost and full-cost references on six of nine revived datasets) once a suitable predictor is deployed.
\end{enumerate}

In summary, the proposed methodology integrates distribution-aware granular information with explicit predictive feasibility constraints. Its theoretical bounds apply strictly to the defined representation and search mechanisms, whereas its predictive efficacy is validated empirically across specified learners, cost structures, and data splits. Ultimately, the framework delivers a cost-efficient attribute subset guaranteed to satisfy internal predictive feasibility; where a simple mean-distribution predictor already matches the full-attribute baseline, the method correctly identifies and returns an empty subset at zero cost, and when a stronger predictor makes a domain informative again, it certifies low-cost subsets that dominate wholesale acquisition.
